(1) 用谐波时间因子 \(e^{i\omega t}\)，线上的电压、流向负载的电流可写 \(V(x)=V_+e^{-i\beta x}+V_-e^{i\beta x}\)、\(I(x)=[V_+e^{-i\beta x}-V_-e^{i\beta x}]/Z_t\)。由负载边界 \(V(\ell)=R_LI(\ell)\) 消去波幅，得到 \(Z_{\rm in}=Z_t\dfrac{R_L+iZ_t\tan(\beta\ell)}{Z_t+iR_L\tan(\beta\ell)}\)，该式为无耗线的一般输入阻抗。
(2) 当 \(\beta\ell=\pi/2\) 时可取 \(\tan(\beta\ell)\to\infty\)，于是 \(Z_{\rm in}=Z_t^2/R_L\)。与馈线阻抗 \(Z_0\) 匹配要求 \(Z_{\rm in}=Z_0\)，所以应选 \(Z_t=\sqrt{Z_0R_L}\)，这是实电阻负载与无损四分之一波长段的理想中心频率匹配条件。
(3) 令 \(\tan(\pi/2+\delta)=-1/\delta+O(\delta)\)，代回阻抗式并展开，\(Z_{\rm in}=\dfrac{Z_t^2}{R_L}\left[1+i\delta\left(\dfrac{R_L}{Z_t}-\dfrac{Z_t}{R_L}\right)+O(\delta^2)\right]\)。在中心匹配条件下 \(Z_t^2/R_L=Z_0\)，由 \(\Gamma_{\rm in}=(Z_{\rm in}-Z_0)/(Z_{\rm in}+Z_0)\) 得 \(\Gamma_{\rm in}\simeq\dfrac{i\delta}{2}(R_L/Z_t-Z_t/R_L)\)。失谐幅度一阶变化，反射功率二阶变化。
(4) 馈线参考面中的入射平均功率为 \(P_+=|V_+|^2/(2Z_0)\)；无耗变换器中能量只被反射或送到负载，所以 \(P_L=P_+(1-|\Gamma_{\rm in}|^2)\)。仅在 \(\delta=0\) 才有 \(\Gamma_{\rm in}=0\)（或恰有 \(R_L=Z_0\) 的无须变换情形）；频率离开设计点后一般出现反射，因此这一结构是窄带匹配器而非理想全频段变压器。
